\documentclass[11pt,a4paper]{article}
\usepackage[T1]{fontenc}
\usepackage{isabelle,isabellesym}
\usepackage{amsmath}
\usepackage{amssymb}
\usepackage{pdfsetup}

\urlstyle{rm}
\isabellestyle{it}
\setlength{\emergencystretch}{3em}

\begin{document}

\title{The F{\"o}ldes--Hammer Characterization of Split Graphs}
\author{Arthur Freitas Ramos \and David Barros Hulak \and Ruy J.~G.~B. de Queiroz}
\maketitle

\begin{abstract}
This entry formalizes the classical F{\"o}ldes--Hammer theorem characterizing
finite split graphs as exactly the graphs with no induced copy of $2K_2$, $C_4$,
or $C_5$.  It builds on the AFP session \texttt{Undirected\_Graph\_Theory}.

\end{abstract}

\section{Overview}

A finite graph is split when its vertex set can be partitioned into a clique and
an independent set. The main theorem of this entry is
\texttt{foldes\_hammer} in \texttt{Foldes\_Hammer.thy}, which proves the
equivalence between split graphs and the absence of induced $2K_2$, $C_4$, and
$C_5$.

The forward direction is established by analyzing how any split partition
excludes each forbidden configuration. The converse direction follows an
extremal argument: choose a maximum clique whose complement minimizes the number
of induced edges, orient any remaining outside edge according to nested
nonneighbor sets in the clique, and swap one clique vertex with an outside
vertex to derive a contradiction.

\section{Sources}

The entry follows the classical theorem due to F{\"o}ldes and
Hammer~\cite{FoldesHammer1977} and standard graph-class background as presented
by Brandst{\"a}dt, Le, and Spinrad~\cite{Brandstadt1999}.

\tableofcontents

\sloppy
\input{session}
\fussy

\bibliographystyle{abbrv}
\bibliography{root}

\end{document}
